\documentclass[11pt]{article}
\usepackage[margin=0.82in]{geometry}
\usepackage{amsmath,amssymb,amsthm,booktabs,graphicx,microtype}
\usepackage[hidelinks]{hyperref}
\usepackage[T1]{fontenc}

\newtheorem{theorem}{Theorem}
\newtheorem{corollary}[theorem]{Corollary}
\newcommand{\R}{\mathbb R}
\newcommand{\Sph}{S^{n-1}}
\newcommand{\vol}{V}

\title{Basis-mixture extremals for the sine transform\\
\large A partial answer to Maresch--Schuster (arXiv:1207.7266)}
\author{Automated mathematical-discovery run \texttt{fa\_banach\_001}}
\date{11 August 2026}

\begin{document}
\maketitle

\noindent\textbf{Status.} Partial result, likely valid, subject to human review.

\section{Source conjecture}

For a finite Borel measure $\mu$ on $\Sph$, its sine transform is
\[
 (\mathcal S\mu)(x)=\int_{\Sph}\sqrt{\lVert x\rVert^2-(x\cdot u)^2}\,d\mu(u).
\]
If $\mu$ is even and isotropic,
$\int u\otimes u\,d\mu(u)=I_n$, this is the support function of a
centrally symmetric convex body denoted $S_\mu$; its polar is
$S_\mu^*$.  For an orthonormal basis $E=(e_1,\dots,e_n)$, write
\[
 \nu_E=\frac12\sum_{i=1}^n(\delta_{e_i}+\delta_{-e_i})
\]
for the associated cross measure.  On page 14, after Theorem 4.3,
Maresch and Schuster conjecture:

\begin{quote}
Among even isotropic measures, $V(S_\mu^*)$ is maximized precisely by
cross measures, while $V(S_\mu)$ is minimized precisely by cross measures.
\end{quote}

\begin{center}
\includegraphics[width=0.91\textwidth]{figures/open_problem_crop.png}
\end{center}

\section{New partial theorem}

Call an even measure an \emph{orthogonal-basis mixture} if there is a
Borel probability measure $\pi$ on $O(n)$ such that
\begin{equation}\label{eq:mixture}
 \mu=\int_{O(n)}R_\#\nu_E\,d\pi(R).
\end{equation}
This is exactly the closed convex hull of the cross measures.  It includes
finite and continuous mixtures (for example normalized spherical Lebesgue
measure), but is a proper subclass of the even isotropic measures.

\begin{theorem}[Both conjectured inequalities on the basis-mixture class]
\label{thm:main}
Let $n\ge3$, let $\nu_E$ be a cross measure, and let $\mu$ have the
representation \eqref{eq:mixture}.  Then
\[
 \vol(S_\mu^*)\le \vol(S_{\nu_E}^*)
 \qquad\text{and}\qquad
 \vol(S_\mu)\ge \vol(S_{\nu_E}).
\]
Equality in either inequality holds if and only if $\mu$ itself is a
cross measure.
\end{theorem}

\section{Proof intuition}

All cross measures are rotations of one model, so all their sine-transform
bodies are rotated copies of one body $K_0$.  Because the sine transform is
linear in the measure, $S_\mu$ is the Minkowski average of these rotated
copies.  Volume to the power $1/n$ grows under Minkowski averaging, which
gives the minimum of the primal volume.  Polar volume has the integral
formula $\frac1n\int h^{-n}$; strict convexity of $t\mapsto t^{-n}$ makes it
shrink under averaging.  The equality cases force almost every rotated body
in the mixture to be the same.  Injectivity of the sine transform on even
measures then forces the measure, not only its body, to be a single cross.

\section{Proof}

Put $K_0=S_{\nu_E}$ and $K_R=RK_0=S_{R_\#\nu_E}$.  Every
$R_\#\nu_E$ is even and isotropic, hence so is $\mu$.  Linearity and
orthogonal equivariance of the sine transform give
\begin{equation}\label{eq:Minkowski}
 h_{S_\mu}(x)=\int_{O(n)}h_{K_R}(x)\,d\pi(R).
\end{equation}
Thus $\overline K:=S_\mu$ is the Minkowski integral of the $K_R$.

For the polar inequality, the polar-coordinate volume formula and Jensen's
inequality yield
\begin{align*}
 \vol(\overline K^*)
 &=\frac1n\int_{\Sph}\left(\int h_{K_R}(u)\,d\pi(R)\right)^{-n}du\\
 &\le \int_{O(n)}\frac1n\int_{\Sph}h_{K_R}(u)^{-n}\,du\,d\pi(R)
 =\vol(K_0^*).
\end{align*}

For the primal inequality, multilinearity and continuity of mixed volume give
\[
 \vol(\overline K)
 =\int_{O(n)}V(\overline K[n-1],K_R)\,d\pi(R).
\]
Minkowski's first inequality, and $\vol(K_R)=\vol(K_0)$, imply
\[
 V(\overline K[n-1],K_R)
 \ge \vol(\overline K)^{(n-1)/n}\vol(K_0)^{1/n}.
\]
After integration and division by the positive common factor,
$\vol(\overline K)^{1/n}\ge\vol(K_0)^{1/n}$.

It remains to check rigidity.  Equality in the polar argument and strict
convexity of $t^{-n}$ imply that $h_{K_R}(u)$ is independent of $R$ for
$\pi$-almost every $R$ and almost every $u$.  Continuity of support functions
then makes the bodies $K_R$ identical almost surely.  In the primal argument,
equality makes the nonnegative gap in Minkowski's first inequality vanish
almost surely.  Hence almost every $K_R$ is homothetic to $\overline K$.
All $K_R$ are centrally symmetric about the origin and have the same volume;
under equality $\overline K$ has that volume as well.  The homotheties
therefore have ratio one and no translation, so again the $K_R$ are identical
almost surely.

Finally, the sine transform is injective on even measures
\cite[Proposition 2.1]{MS}.  Thus equality of $K_R=S_{R_\#\nu_E}$ implies
that the measures $R_\#\nu_E$ are identical almost surely.  Equation
\eqref{eq:mixture} then says that $\mu$ is that cross measure.  The converse
is immediate. \qed

\section{A discrete criterion}

Let $\mu$ be supported on finitely many antipodal lines $L_1,\dots,L_N$,
and let $a_j$ be its total mass on $L_j$.  Form the hypergraph whose edges
are the $n$-element subsets $B$ of pairwise orthogonal lines (hence
orthonormal bases after choosing unit representatives).

\begin{corollary}[Fractional orthogonal-basis decomposition]
If there are numbers $\lambda_B\ge0$ such that
\[
 \sum_B\lambda_B=1,
 \qquad
 a_j=\sum_{B\ni L_j}\lambda_B\quad(1\le j\le N),
\]
then both inequalities and both equality conclusions of Theorem
\ref{thm:main} hold for $\mu$.
\end{corollary}

\begin{proof}
The displayed equations say exactly that $\mu$ is the convex combination,
with coefficients $\lambda_B$, of the cross measures carried by the basis
edges $B$.  Apply Theorem \ref{thm:main}.
\end{proof}

\section{Scope, upgrade attempts, and novelty}

The theorem does not settle the full conjecture.  A regular-simplex tight
frame is isotropic but its support contains no orthogonal basis, so it cannot
belong to the basis-mixture class.  Moreover, that class is closed: it is the
continuous image of the compact space of probability measures on $O(n)$.
Thus a density argument cannot remove the obstruction.

The attack also tested the most immediate counterexample route.  In dimension
three, deterministic spherical sampling put the polar volume of the regular
tetrahedral frame below that of the cross, and random searches over Parseval
frames with four through nine lines found no violation.  A separate zonoid
determinant computation put the tetrahedral primal volume above the cross
value.  These computations are evidence only; the proof above does not use
them.

A bounded literature search through 11 August 2026 found no full resolution.
The complex analogue still described both the real and complex problems as
open in 2015 \cite{HH}.  Li--Xi--Zhang's Grassmannian theorem proves sharp
reverse inequalities when the Grassmann rank divides the dimension and
explicitly leaves the non-divisible case open \cite[Theorem 1.2]{LXZ}.
The sine transform has rank $n-1$, which does not divide $n$ for $n\ge3$.
Exact-title, exact-conjecture, author/citation, and
``basis/mixture of cross measures'' searches found no prior version of the
partial theorem.  Novelty is therefore plausible but provisional.

\begin{thebibliography}{9}
\bibitem{MS}
G. Maresch and F. E. Schuster,
\emph{The sine transform of isotropic measures},
Int. Math. Res. Not. IMRN 2012 (2012), 717--739;
arXiv:1207.7266.

\bibitem{HH}
Q. Huang and B. He,
\emph{Volume inequalities for complex isotropic measures},
Geom. Dedicata 177 (2015), 401--428.

\bibitem{LXZ}
A.-J. Li, D. Xi, and G. Zhang,
\emph{Volume inequalities of convex bodies from cosine transforms on
Grassmann manifolds}, Adv. Math. 304 (2017), 494--538.
\end{thebibliography}

\end{document}
